\originalchapter{来源对照}
\noindent 下表按实际下载的 PDF 正文建立对应, 页数为源文件物理页数. 官网主题表有错位: Cantor 集与非 Borel 可测集为第3章编者补充; 原第17讲没有平滑稠密证明, 该内容实际在第21讲. 原第19讲已经陈述一般 Fubini, 第20讲已经给出 Young 不等式. 原第24讲只证明最大算子的专门插值估计. 中文补证、勘误与新增练习见各章正文及独立核查记录.

\begin{longtable}{@{}p{.07\textwidth}p{.07\textwidth}p{.66\textwidth}p{.11\textwidth}@{}}
\toprule 讲 & 源页 & 实际内容 & 本册章\\\midrule\endhead
\href{https://ocw.mit.edu/courses/18-125-measure-and-integration-fall-2003/resources/18125_lec1/}{1} & 3 & 可测空间、可测复合与Borel集 & 1 \\
\href{https://ocw.mit.edu/courses/18-125-measure-and-integration-fall-2003/resources/18125_lec2/}{2} & 4 & 实可测函数、极限、简单函数与非负积分 & 1,2 \\
\href{https://ocw.mit.edu/courses/18-125-measure-and-integration-fall-2003/resources/18125_lec3/}{3} & 3 & Riemann积分判据、测度连续性、积分性质 & 1,2 \\
\href{https://ocw.mit.edu/courses/18-125-measure-and-integration-fall-2003/resources/18125_lec4/}{4} & 4 & 单调收敛、积分可加性、级数与Fatou & 2 \\
\href{https://ocw.mit.edu/courses/18-125-measure-and-integration-fall-2003/resources/18125_lec5/}{5} & 5 & 复积分、主导收敛、零集和完备化 & 2 \\
\href{https://ocw.mit.edu/courses/18-125-measure-and-integration-fall-2003/resources/18125_lec6/}{6} & 4 & 矩形与多面体、开紧集、内外测度 & 3 \\
\href{https://ocw.mit.edu/courses/18-125-measure-and-integration-fall-2003/resources/18125_lec7/}{7} & 4 & 可测集定义和Lebesgue测度可列可加性 & 3 \\
\href{https://ocw.mit.edu/courses/18-125-measure-and-integration-fall-2003/resources/18125_lec8/}{8} & 3 & Carathéodory准则、开集紧集逼近、完备性 & 3 \\
\href{https://ocw.mit.edu/courses/18-125-measure-and-integration-fall-2003/resources/18125_lec9/}{9} & 4 & 不变性、Vitali不可测集、线性变换 & 3 \\
\href{https://ocw.mit.edu/courses/18-125-measure-and-integration-fall-2003/resources/18125_lec10/}{10} & 5 & Riesz表示、由Riemann泛函构造Lebesgue测度 & 4 \\
\href{https://ocw.mit.edu/courses/18-125-measure-and-integration-fall-2003/resources/18125_lec11/}{11} & 3 & Lusin、Vitali–Carathéodory & 4 \\
\href{https://ocw.mit.edu/courses/18-125-measure-and-integration-fall-2003/resources/18125_lec12/}{12} & 3 & 连续函数逼近、几乎处处收敛和集合可加积分 & 2,4 \\
\href{https://ocw.mit.edu/courses/18-125-measure-and-integration-fall-2003/resources/18125_lec13/}{13} & 4 & Egorov、依测度收敛、子列及主导收敛 & 4 \\
\href{https://ocw.mit.edu/courses/18-125-measure-and-integration-fall-2003/resources/18125_lec14/}{14} & 4 & 凸函数、Jensen、Hölder、Minkowski & 5 \\
\href{https://ocw.mit.edu/courses/18-125-measure-and-integration-fall-2003/resources/18125_lec15/}{15} & 5 & Lp商空间、范数与Riesz–Fischer完备性 & 5 \\
\href{https://ocw.mit.edu/courses/18-125-measure-and-integration-fall-2003/resources/18125_lec16/}{16} & 2 & Cc在Lp稠密、Cc在C0稠密 & 5 \\
\href{https://ocw.mit.edu/courses/18-125-measure-and-integration-fall-2003/resources/18125_lec17/}{17} & 4 & Lp与lp包含、局部Lp、凸性 & 5 \\
\href{https://ocw.mit.edu/courses/18-125-measure-and-integration-fall-2003/resources/18125_lec18/}{18} & 6 & Rn非负及可积Fubini、截面论证和双积分例 & 6 \\
\href{https://ocw.mit.edu/courses/18-125-measure-and-integration-fall-2003/resources/18125_lec19/}{19} & 3 & 可测直积、一般乘积测度与Fubini & 6 \\
\href{https://ocw.mit.edu/courses/18-125-measure-and-integration-fall-2003/resources/18125_lec20/}{20} & 3 & 卷积、Young、共轭卷积一致连续 & 7 \\
\href{https://ocw.mit.edu/courses/18-125-measure-and-integration-fall-2003/resources/18125_lec21/}{21} & 4 & 近似恒等、磨光、平滑稠密 & 7 \\
\href{https://ocw.mit.edu/courses/18-125-measure-and-integration-fall-2003/resources/18125_lec22/}{22} & 5 & 绝对连续、Vitali覆盖与弱L1最大估计 & 8,9 \\
\href{https://ocw.mit.edu/courses/18-125-measure-and-integration-fall-2003/resources/18125_lec23/}{23} & 5 & Lebesgue微分、Lebesgue点、积分函数求导 & 8 \\
\href{https://ocw.mit.edu/courses/18-125-measure-and-integration-fall-2003/resources/18125_lec24/}{24} & 5 & 广义Minkowski、分布函数、最大算子Lp估计 & 8 \\
\bottomrule
\end{longtable}
第9章的积分密度与 Lebesgue 分解是编者补充, 不对应上述任何源 PDF; 其中绝对连续函数一节补证第22讲只陈述的逆向积分定理. 课程大纲中的 Hausdorff 测度、面积与余面积公式没有被包含在实际讲义中, 因此尚未译编. 课程描述提及 Fourier transform, 这24份文件也没有对应正文. 独立作业集与外部教材原题不属于本册的来源范围.
